% Báo cáo rà soát lỗi Hệ thống Thời khóa biểu — 30/8/2026
% Biên dịch: pdflatex (cần gói babel vietnamese) hoặc xelatex đều được.
\documentclass[12pt,a4paper]{article}
\usepackage[utf8]{inputenc}
\usepackage[T5]{fontenc}
\usepackage[vietnamese]{babel}
\usepackage[margin=2.2cm]{geometry}
\usepackage{longtable}
\usepackage{booktabs}
\usepackage{enumitem}
\usepackage{xcolor}
\usepackage[colorlinks=true,linkcolor=blue!60!black]{hyperref}

\newcommand{\muc}[1]{\textbf{#1}}
\definecolor{do}{RGB}{180,30,30}
\definecolor{xanh}{RGB}{21,128,61}
\newcommand{\nghiemtrong}{\textcolor{do}{\textbf{Nghiêm trọng}}}
\newcommand{\cao}{\textcolor{do}{Cao}}
\newcommand{\vua}{Vừa}
\newcommand{\nhe}{Nhẹ}
\newcommand{\dasua}{\textcolor{xanh}{\textbf{Đã sửa}}}

\title{\textbf{BÁO CÁO RÀ SOÁT LỖI}\\[4pt]
\large Hệ thống Thời khóa biểu --- trường tiểu học nhiều phân hiệu}
\author{Rà soát tự động bằng 7 bộ kiểm thử + 3 tác tử soi mã độc lập}
\date{Ngày 30 tháng 8 năm 2026}

\begin{document}
\maketitle

\section*{Tóm tắt}
Đợt rà soát tiến hành theo hai hướng song song: (1) chạy \textbf{toàn bộ bảy bộ
kiểm thử tự động} của dự án; (2) giao \textbf{ba tác tử soi mã độc lập} rà ba
vùng rủi ro cao nhất theo lịch sử sự cố --- tầng truy cập dữ liệu (DULIEU),
thuật toán và nghiệp vụ (LOGIC), giao diện và bảo mật phía trình duyệt.

Kết quả: bảy bộ kiểm thử \textbf{đều xanh} trước khi rà (1.524 phép thử), nhưng
ba tác tử tìm ra \textbf{13 lỗi thật} chưa bộ nào canh tới --- trong đó 2 lỗi
nghiêm trọng, 1 lỗi mức cao, 1 lỗi bảo mật XSS. \textbf{Cả 13 lỗi đã được sửa
trong cùng ngày}, kèm 9 phép thử mới canh không cho tái phát. Sau khi vá, cả
bảy bộ kiểm thử chạy lại đều xanh (1.533 phép thử).

\section{Kết quả bảy bộ kiểm thử (trước khi vá)}
\begin{longtable}{@{}llr@{}}
\toprule
\muc{Bộ} & \muc{Phạm vi} & \muc{Kết quả}\\
\midrule
\texttt{npm test}          & hàm thuần (thuật toán, dữ liệu, quyền)      & 455 đạt, 0 hỏng\\
\texttt{npm run soi}       & giao diện 12+ màn hình, trình duyệt giả      & 488 đạt, 0 hỏng\\
\texttt{npm run soi-mau}   & hình thức tệp Excel xuất ra                  & 81 đạt, 0 hỏng\\
\texttt{npm run soi-nhap}  & đường nhập liệu của trường mới tinh          & 60 đạt, 0 hỏng\\
\texttt{npm run soat}      & cú pháp SQL, quy tắc RLS, cột quyền          & 409 câu lệnh, 0 lỗi\\
\texttt{soi-worker}        & Web Worker, CSP, nút Xếp kỹ                  & 9 đạt, 0 hỏng\\
\texttt{npm run kiemdinh}  & trình kiểm định thuật toán độc lập, 7 kịch bản & 22 đạt, 0 hỏng\\
\bottomrule
\end{longtable}

Nhận xét: các bộ kiểm thử canh rất chắc những gì chúng \emph{đã biết}; 13 lỗi
dưới đây đều nằm ở những đường đi mà chưa phép thử nào bước qua.

\section{Bảng tổng hợp 13 lỗi}
\begin{longtable}{@{}p{0.6cm}p{1.9cm}p{9.3cm}p{1.6cm}@{}}
\toprule
\muc{\#} & \muc{Mức} & \muc{Lỗi} & \muc{T.thái}\\
\midrule
1 & \nghiemtrong & Nhập Excel ma trận \textbf{xoá sạch Gmail, điện thoại, ghi chú,
    phân hiệu} và đè định mức riêng của mọi giáo viên đang có trên máy chủ & \dasua\\
2 & \nghiemtrong & Vân tay ``chưa lưu'' \textbf{không phủ khung giờ và các cột mới}
    (Gmail, điện thoại, số tiết chuẩn môn) --- sửa xong tải lại trang là mất âm thầm & \dasua\\
3 & \cao & Chỉnh tay TKB: chiều ngược của phép hoán đổi \textbf{thiếu kiểm buổi bận}
    --- vi phạm ràng buộc cứng số 7, tiết bị ghim nên Xếp lại cũng không chữa & \dasua\\
4 & \cao & \textbf{XSS} qua tên lớp tại 3 chỗ chèn thẳng \texttt{l.ten} không qua
    \texttt{esc()} (Bảng điều hành, mục Phân hiệu, hộp Phân bổ lớp) & \dasua\\
5 & \vua & Bảng \textbf{ngày công} đếm đơn nghỉ chồng lấn (buổi + cả ngày cùng
    một ngày) thành 1,5 công thay vì 1 công --- sai số liệu báo cáo hằng tháng & \dasua\\
6 & \vua & Chốt chặn cuối trước khi lưu \textbf{dạy thay} yếu hơn bộ lọc ứng viên:
    không bắt ``cùng buổi, khác tiết, hai phân hiệu'' và không kiểm giới hạn tiết/buổi & \dasua\\
7 & \vua & Đổi \textbf{tên phân hiệu} bị ghi thành tạo-mới-xoá-cũ: giáo viên mất
    nhãn phân hiệu, PHT phụ trách mất phạm vi lưu mà máy vẫn báo ``Đã lưu'' & \dasua\\
8 & \vua & Chế độ chủ hệ thống xem trường khác: tải lỗi thì \textbf{bày dữ liệu mẫu}
    dưới thẻ ghi tên trường thật, dễ đọc nhầm thành dữ liệu của trường ấy & \dasua\\
9 & \vua & \textbf{Đăng xuất không dọn} cờ trường-đang-xem: tài khoản đăng nhập sau
    trên cùng máy bị tải nhầm trường, thấy ``trường chưa có dữ liệu'' & \dasua\\
10 & \vua & Danh mục môn/phòng: xoá xong mà ghi lại hỏng thì \textbf{nuốt lỗi} ---
    máy chủ mất sạch danh mục nhưng vẫn báo ``Đã nhập'' & \dasua\\
11 & \vua & Mở app lúc \textbf{sóng yếu} bị xoá vé đăng nhập còn tốt --- hôm sau
    thầy cô lại phải gõ mật khẩu; vé mới xoay vòng chưa kịp ghi cũng bị vứt & \dasua\\
12 & \nhe & Bốn đường ghi trong-năm-học (buổi bận, dạy thay, báo nghỉ, đánh dấu
    xử lý) \textbf{thiếu chốt chỉ-đọc} của chế độ xem trường khác & \dasua\\
13 & \nhe & Id vừa khai trong app (chưa lưu) lọt vào cột \texttt{uuid} ở mã mời,
    dạy thay, báo nghỉ --- người dùng nhận câu lỗi tiếng Anh khó hiểu & \dasua\\
\bottomrule
\end{longtable}

\section{Chi tiết các lỗi chính và cách sửa}

\subsection*{Lỗi 1 --- Nhập ma trận xoá Gmail cả trường (nghiêm trọng)}
\muc{Kịch bản:} mẫu ma trận (mẫu \emph{duy nhất} được mời ở màn Phân công) không
có cột Gmail, điện thoại, định mức, phân hiệu. Trường đã khai Gmail cho 35 thầy
cô, giữa kỳ tải mẫu về sửa phân công rồi nhập lại: dòng giáo viên upsert theo mã
\emph{khớp đúng dòng cũ} trên máy chủ và ghi đè các cột ấy bằng \texttt{null} ---
sáng hôm sau cả trường không ai đăng nhập bằng Google được nữa; định mức riêng
của giáo viên kiêm nhiệm (18 tiết) bị nâng lại 23 nên quy tắc R01 im lặng.

\muc{Cách sửa:} tại đường ghi duy nhất \texttt{ghiDuLieuNguon()}: ô để trống thì
\emph{giữ nguyên giá trị cũ} đang có (chỉ tin dữ liệu cũ khi nó thật sự đến từ
máy chủ); trình đọc Excel đánh dấu \texttt{dinhMucKhai} để phân biệt ``ô bỏ
trống'' với ``người dùng khai 23''. Đúng nguyên tắc đã ghi trong CLAUDE.md:
\emph{``Ô Gmail để trống thì giữ nguyên địa chỉ cũ, không xoá.''}

\subsection*{Lỗi 2 --- Vân tay ``chưa lưu'' thủng ở các cột mới (nghiêm trọng)}
\muc{Kịch bản:} cơ chế \texttt{vanTayNguon()} sinh ra sau sự cố mất dữ liệu
2/8/2026, nhưng không được cập nhật khi bảng giáo viên thêm cột Gmail, điện
thoại, ghi chú, phân hiệu chính, mã GV; khung giờ và số tiết chuẩn của môn cũng
không nằm trong vân tay. Quản trị gõ Gmail cho 35 thầy cô ngay trong bảng, quên
bấm Lưu, trang tải lại --- mất sạch mà không dải đỏ, không \texttt{beforeunload}.

\muc{Cách sửa:} bổ sung đủ các trường trên vào vân tay; thêm cụm phép thử 18b2
sửa từng trường một và đòi \texttt{coThayDoiChuaLuu()} phải đổi.

\subsection*{Lỗi 3 --- Hoán đổi tay đẩy tiết vào buổi đồng nghiệp đã báo bận (cao)}
\muc{Kịch bản:} \texttt{kiemTraChuyen()} kiểm buổi bận cho chiều xuôi nhưng
\emph{quên chiều ngược}: kéo tiết Toán thả vào ô Tiếng Anh của cô Mận thì tiết
của cô bị đẩy vào đúng buổi cô đã đăng ký bận --- vi phạm ràng buộc cứng số 7
một cách âm thầm, và vì tiết chỉnh tay được ghim nên lần Xếp lại sau cũng không
chữa. Chốt của đường chỉnh tay yếu hơn chốt của máy (\texttt{doiChoDuoc()} kiểm
đủ cả hai chiều).

\muc{Cách sửa:} thêm phép kiểm \texttt{S.gvNghi} cho chiều ngược, câu báo nói rõ
ai và vì sao.

\subsection*{Lỗi 4 --- XSS qua tên lớp (bảo mật)}
\muc{Kịch bản:} ba chỗ chèn thẳng tên lớp vào HTML không qua \texttt{esc()}.
Tên lớp là chuỗi nhập tự do (sửa tại chỗ hoặc cột \texttt{Ten\_lop} của Excel)
và CSP còn \texttt{'unsafe-inline'} (giới hạn đã ghi nhận của quy ước một tệp),
nên một tên lớp chứa mã có thể chạy trong phiên của bất kỳ ai mở trang --- kể cả
chủ hệ thống đang xem trường khác. Vé đăng nhập nằm ở \texttt{localStorage}.

\muc{Cách sửa:} bọc \texttt{esc()} cả ba chỗ, đồng nhất với mọi nơi khác.

\subsection*{Lỗi 5 --- Ngày công đếm trùng đơn chồng lấn}
Sáng gửi ``nghỉ buổi sáng'', trưa gửi tiếp ``nghỉ cả ngày'' là hai dòng cho cùng
một ngày; bảng tổng hợp cộng thành 1,5 công. Nay quy công theo \emph{ngày}: có
``cả ngày'' thì buổi lẻ cùng ngày không đếm nữa, một ngày tối đa 1 công.

\subsection*{Lỗi 6 --- Chốt cuối dạy thay yếu hơn bộ lọc}
\texttt{xungDotDayThay()} (chạy ngay trước khi lưu) không cộng các tiết dạy thay
đã nhận vào tập phân hiệu trong buổi và không kiểm giới hạn tiết/buổi --- hai
quản lý ở hai máy cùng chọn một người cho hai phân hiệu (khác tiết, cùng buổi)
đều thấy ``sạch''. Nay chốt cuối kiểm đủ cả dòng đã lưu lẫn dòng trong chính mẻ
đang lưu. \emph{Lưu ý còn để ngỏ:} hàng rào tuyệt đối cho kịch bản hai máy phải
nằm ở cơ sở dữ liệu (xem mục~\ref{sec:dengo}).

\subsection*{Lỗi 7 --- Đổi tên phân hiệu làm đứt tham chiếu}
Đổi ``Điểm trường Diễn Liên'' thành ``Phân hiệu Diễn Liên'' rồi bấm Lưu: bản cũ
dò phân hiệu theo tên nên tạo dòng mới rồi xoá dòng cũ --- giáo viên của phân
hiệu ấy mất nhãn (\texttt{on delete set null}), tài khoản PHT phụ trách trỏ vào
id đã xoá nên phạm vi lưu rỗng: bấm Lưu được báo thành công mà không ghi lớp
nào. Nay dòng đã có trên máy chủ thì \emph{sửa tên đúng dòng ấy} qua
\texttt{suaHang()}; sửa không được thì lùi về đường cũ, không hỏng thêm gì.

\subsection*{Các lỗi 8--13}
Đã sửa theo đúng khuôn sẵn có của dự án: điều kiện khôi phục trong
\texttt{moTruongDeXem()} xét cả \texttt{ok=false} và nguồn dữ liệu mẫu;
\texttt{dangXuat()} dọn \texttt{truongXem}/\texttt{vanBan}/\texttt{banCongBo};
danh mục môn/phòng chỉ nuốt lỗi ``bảng chưa có'', còn lại báo ra kèm lời nhắc
bấm Lưu lại; vé đăng nhập chỉ bị xoá khi máy chủ từ chối (4xx), lỗi mạng thì giữ
vé và ghi vé xoay vòng xuống máy \emph{ngay} sau khi làm mới; sáu đường ghi
trong-năm-học thêm chốt \texttt{dangXemTruongKhac()}; lỗi
\texttt{invalid input syntax for type uuid} được dịch thành việc phải làm
(``Bấm Lưu ở mục Giáo viên trước\dots'') thay vì câu tiếng Anh.

\section{Phép thử mới (9 phép, đều đã thử ngược)}
\begin{itemize}[leftmargin=1.4em,itemsep=1pt]
  \item Nhập tệp không có cột Gmail giữ nguyên Gmail, điện thoại, ghi chú, định mức riêng (mục 11).
  \item \texttt{duLieuTuBang}: ô \texttt{Dinh\_muc} bỏ trống đánh dấu \texttt{dinhMucKhai = 0} (mục 11).
  \item Cùng buổi, khác tiết, hai phân hiệu --- chốt cuối dạy thay phải bắt, kể cả hai dòng cùng mẻ (mục 14).
  \item Đơn chồng lấn buổi + cả ngày chỉ tính 1 công (mục 18b).
  \item Cụm 18b2 --- năm phép: sửa Gmail, điện thoại/ghi chú/phân hiệu/mã GV, số tiết khung giờ,
        bật tắt buổi, số tiết chuẩn/màu môn đều làm vân tay đổi.
\end{itemize}

\section{Kết quả sau khi vá}
Cả bảy bộ chạy lại xanh trọn: \texttt{npm test} 464 đạt (455 cũ + 9 mới),
\texttt{soi} 488, \texttt{soi-mau} 81, \texttt{soi-nhap} 60, \texttt{soat} 0
lỗi, \texttt{soi-worker} 9, \texttt{kiemdinh} 22 --- tổng \textbf{1.533 phép
thử, 0 hỏng}. Các mốc thuật toán giữ nguyên: 710/710 tiết một phân hiệu,
1698/1698 ba phân hiệu, 0 xung đột.

\section{Việc còn để ngỏ (khuyến nghị, chưa sửa trong đợt này)}\label{sec:dengo}
\begin{enumerate}[leftmargin=1.6em,itemsep=2pt]
  \item \muc{Trigger ở cơ sở dữ liệu cho dạy thay hai máy.} Chốt cuối phía app đã
    vá, nhưng hai trình duyệt chỉ biết trạng thái của mình; hàng rào tuyệt đối
    cho ``một giáo viên, một buổi, một phân hiệu'' trên bảng \texttt{day\_thay}
    cần một trigger \texttt{before insert} (chỉ số unique không diễn tả được
    điều kiện này). Nên làm thành một tệp \texttt{db/*.sql} riêng, thử kỹ rồi
    chạy trên máy chủ thật.
  \item \muc{Chặn chồng lấn ngay lúc gửi báo nghỉ.} Hiện mới sửa ở bảng tổng hợp;
    gọn hơn nữa là khi gửi ``cả ngày'' thì thay các dòng buổi lẻ cùng ngày.
  \item \muc{Tham số \texttt{boQua} chưa được nối dây} ở ba nơi gọi thật của bộ
    dạy thay --- hiện không sai vì chưa có luồng ``sửa phương án đã lưu'', nhưng
    ai thêm luồng ấy phải nhớ truyền.
  \item \muc{\texttt{esc()} không thoát dấu nháy đơn} --- an toàn với kỷ luật hiện
    tại (mọi thuộc tính động dùng nháy kép); giữ kỷ luật ấy hoặc bổ sung khi
    thuận tiện.
\end{enumerate}

\bigskip
\noindent\textit{Toàn bộ thay đổi nằm trong \texttt{src/index.html} và
\texttt{test/kiem-thu.mjs}, không đụng lược đồ cơ sở dữ liệu --- không cần chạy
tệp SQL nào trên máy chủ cho đợt vá này.}

\end{document}
